\documentclass[11pt,a4paper]{article}

\usepackage[utf8]{inputenc}
\usepackage[T1]{fontenc}
\usepackage[brazil]{babel}
\usepackage{lmodern}
\usepackage{textcomp}
\usepackage[a4paper,margin=2.5cm]{geometry}
\usepackage{amsmath,amssymb,amsthm}
\usepackage{booktabs}
\usepackage{graphicx}
\usepackage{float}
\usepackage{longtable}
\usepackage{array}
\newcolumntype{L}[1]{>{\raggedright\arraybackslash}p{#1}}
\usepackage[table]{xcolor}
\usepackage{listings}
\usepackage{titlesec}

% --- Paleta de cores do documento (uso discreto e pontual) -----------
\definecolor{accentblue}{RGB}{70,130,180}  % azul discreto, usado apenas em um fio sob \section e em links externos
\definecolor{linkblue}{RGB}{60,110,170}    % azul claro, reservado a links/URLs reais (não à navegação interna)
\definecolor{tablehead}{RGB}{242,246,250}  % cinza-azulado muito claro, fundo de cabeçalho de tabela

% Links de navegação interna (sumário, referências cruzadas, citações)
% permanecem pretos, como em um relatório impresso; apenas URLs/arquivos
% externos usam o azul claro, conforme solicitado.
\usepackage[
	colorlinks=true,
	linkcolor=black,
	citecolor=black,
	urlcolor=linkblue,
	filecolor=linkblue,
	linktoc=all
]{hyperref}
\usepackage{url}
% URLs: além da cor azul, ganham sublinhado. Cada segmento entre "/" é
% sublinhado individualmente com \underline e separado por um ponto de
% quebra (\allowbreak), de modo que o sublinhado sobrevive à quebra de
% linha sem depender de pacotes externos frágeis (soul/ulem mostraram
% corromper o texto quando combinados com \href nesta instalação).
% \vphantom{Ag} força a mesma altura/profundidade de referência em
% todo segmento (independente de ter ou não descendentes como "g"),
% o que mantém o sublinhado na mesma altura em toda a extensão --
% sem isso, cada segmento desenha a linha em uma altura ligeiramente
% diferente, dando a impressão de um traço quebrado/desalinhado.
\newcommand{\segurl}[1]{\underline{\vphantom{Ag}#1}\hspace{0pt}}
\newcommand{\urlPlatformInfo}{\href{https://github.com/FNakano/IAA-Apostila/tree/main/getPlatformInfo}{\ttfamily\color{linkblue}\segurl{https://github.com/}\allowbreak\segurl{FNakano/}\allowbreak\segurl{IAA-Apostila/}\allowbreak\segurl{tree/}\allowbreak\segurl{main/}\allowbreak\segurl{getPlatformInfo}}}
\newcommand{\urlRepoEP}{\href{https://github.com/FNakano/IAA-Apostila/tree/main/EP2}{\ttfamily\color{linkblue}\segurl{https://github.com/}\allowbreak\segurl{FNakano/}\allowbreak\segurl{IAA-Apostila/}\allowbreak\segurl{tree/}\allowbreak\segurl{main/}\allowbreak\segurl{EP2}}}

% Títulos de seção: tipografia preta tradicional, com um fio fino em
% tom azul discreto como único toque de cor -- hierarquia visual sem
% "pintar" o texto do título.
\titleformat{\section}
	{\normalfont\Large\bfseries}{\thesection}{0.8em}{}
	[{\color{accentblue}\titlerule[0.6pt]}]
\titleformat{\subsection}{\normalfont\large\bfseries}{\thesubsection}{0.8em}{}
\titleformat{\subsubsection}{\normalfont\normalsize\bfseries\itshape}{\thesubsubsection}{0.8em}{}

\lstset{
  basicstyle=\ttfamily\small,
  keywordstyle=\color{blue!60!black}\bfseries,
  commentstyle=\color{green!40!black}\itshape,
  stringstyle=\color{red!60!black},
  frame=single,
  breaklines=true,
  columns=fullflexible,
  keepspaces=true,
  showstringspaces=false,
  literate={á}{{\'a}}1 {é}{{\'e}}1 {í}{{\'i}}1 {ó}{{\'o}}1 {ú}{{\'u}}1
           {ã}{{\~a}}1 {õ}{{\~o}}1 {ç}{{\c c}}1 {â}{{\^a}}1 {ê}{{\^e}}1 {ô}{{\^o}}1 {à}{{\`a}}1
           {Á}{{\'A}}1 {É}{{\'E}}1 {Ç}{{\c C}}1 {Ã}{{\~A}}1
}

% Quebra de linha em nomes de arquivo: permite quebrar depois de "_" (e, no
% texto, depois de "/" via \barra) dentro de \texttt, sem inserir hífen, e dá
% folga extra ao espaçamento para evitar linhas que ultrapassem a margem.
\newcommand{\barra}{/\allowbreak}
\let\underscoreoriginal\_
\renewcommand{\_}{\underscoreoriginal\allowbreak}
\setlength{\emergencystretch}{3em}

\newtheorem*{teorema}{Teorema}

\newcommand{\Q}{\mathbb{Q}}
\newcommand{\N}{\mathbb{N}}

% Inclui a figura apenas se o arquivo existir
% #1 = arquivo, #2 = legenda, #3 = rótulo, #4 = fonte (só na legenda)
\newcommand{\figuraopcional}[4]{%
  \IfFileExists{#1}{%
    \begin{figure}[H]
      \centering
      \includegraphics[width=0.82\textwidth]{#1}
      \caption[#2]{#2 #4}
      \label{#3}
    \end{figure}%
  }{%
    \begin{center}\fbox{\parbox{0.8\textwidth}{\centering\small
      Figura pendente: \texttt{\detokenize{#1}}\\ (#2)}}\end{center}%
  }%
}

\begin{document}

% ======================================================================
% CAPA
% ======================================================================
\begin{titlepage}
	\centering
	\vspace*{1.5cm}
	{\large UNIVERSIDADE DE SÃO PAULO\par}
	{\large ESCOLA DE ARTES, CIÊNCIAS E HUMANIDADES\par}
	\vspace{0.2cm}
	{\normalsize Bacharelado em Sistemas de Informação\par}

	\vfill

	{\Huge\bfseries Comparação de Algoritmos de Ordenação\par}
	\vspace{0.35cm}
	\textcolor{accentblue}{\rule{0.4\linewidth}{0.6pt}}
	\vspace{0.35cm}
	{\Large\itshape Quicksort vs.\ Radix Sort\par}
	\vspace{0.6cm}
	{\large EP2 --- Introdução à Análise de Algoritmos\par}

	\vfill

	{\large\bfseries Igor Fernando Casita Ferreira da Silva\par}
	{\normalsize RA 18245598\par}

	\vspace{1.2cm}

	{\normalsize Docente responsável: Prof.\ Dr.\ Fabio Nakano\par}

	\vfill

	{\normalsize São Paulo\par}
	{\normalsize Outubro de 2026\par}
	\vspace*{1cm}
\end{titlepage}

\tableofcontents
\listoffigures
\listoftables
\newpage

%======================================================================
\setcounter{section}{-1}
\section{Contexto, Definições e Objetivo}
%======================================================================

\subsection{Contexto}

Este Relatório documenta o EP2 (Exercício-Programa~2) da disciplina de
Introdução à Análise de Algoritmos (IAA), cujo enunciado pede a
implementação e comparação empírica de dois algoritmos de ordenação,
justificando os resultados observados com base na teoria de análise de
algoritmos estudada na disciplina.

\subsection{Definições}

\begin{description}
  \item[Trabalho:] o conjunto completo de arquivos entregues para este
    EP2: o Relatório, os programas-fonte, os resultados brutos e
    agregados, os gráficos e os scripts de análise estatística.
  \item[Relatório:] o presente documento (\texttt{relatorio.pdf}, gerado
    a partir de \texttt{relatorio.tex}).
  \item[Autor:] Igor Fernando Casita Ferreira da Silva, RA~18245598 ---
    EP2 é trabalho individual.
  \item[Professor:] Prof.\ Dr.\ Fábio Nakano, docente da disciplina de
    Introdução à Análise de Algoritmos.
  \item[Enunciado:] especifica a implementação de dois algoritmos de
    ordenação (escolha livre, desde que sustentem a discussão pedida na
    Questão~5), a execução para vários tamanhos de entrada, o
    preenchimento de uma planilha compartilhada com os tempos medidos, a
    coleta de informações de hardware/software do ambiente de execução
    (conforme \urlPlatformInfo),
    e a geração de uma chave secreta (mesmo método de
    \texttt{calcSecret.c}) para identificar os experimentos na planilha.
    O código-base/esqueleto fornecido está em
    \urlRepoEP.
\end{description}

\subsection{Objetivo}

Implementar \textbf{Quicksort} e \textbf{Radix Sort}, medir seus tempos
de execução reais para tamanhos de entrada de $10^3$ a $2\times10^8$
(e, nos tamanhos oficiais da planilha, até $2\times10^9$), e discutir os
resultados com base na complexidade assintótica (temporal e espacial) de
cada algoritmo, incluindo os casos em que o comportamento observado na
prática diverge do previsto pela análise assintótica de tempo isolada
(Seção~\ref{sec:resultados}).

\textbf{Todas as medições deste Relatório foram feitas no computador
pessoal do autor} (não em ambiente de nuvem ou máquina compartilhada),
especificado na Seção~\ref{sec:ambiente}, usando \texttt{clock()} e
\texttt{gettimeofday()}, conforme especificado no enunciado.

\subsection{Organização dos Arquivos}

\begin{lstlisting}
EP2/
|-- relatorio.pdf / relatorio.tex   # este Relatorio
|-- respostas_EP2_completo.docx/.md # versao docx/markdown (equivalente)
|-- respostas_EP2_resumido.docx/.md # versao resumida
|-- src/
|   |-- rng.c/.h, utilarray.c/.h    # fornecidos pelo professor (sem alteracoes)
|   |-- elapsedRunningTime.c/.h     # fornecido pelo professor (sem alteracoes)
|   |-- quicksort.c/.h              # implementacao propria
|   |-- radixsort.c/.h              # implementacao propria
|   |-- gen_input.c, run_sort.c     # geracao de entrada e execucao medida
|   |-- benchmark_clock.c           # bateria dos tamanhos oficiais da planilha
|   `-- secretkey_gen.c             # chave secreta (nao-interativo)
`-- resultados/
    |-- raw_results.csv             # medicoes repetidas (clock + wall)
    |-- tempos_planilha_clock.csv   # tamanhos oficiais da planilha (so clock)
    |-- resumo_estatistico.csv, regressao.txt
    `-- graficos_r/                 # graficos PNG + analise.R
\end{lstlisting}

%======================================================================
\section{Fundamentação Teórica}
%======================================================================

\subsection{Quicksort}

\textbf{Classe:} ordenação por comparação, \emph{in-place}, por divisão e
conquista (\emph{divide-and-conquer}), \textbf{instável}.

\textbf{Ideia central:} escolhe-se um pivô; particiona-se o array em
elementos menores e maiores que o pivô; ordena-se recursivamente cada
partição.

\subsubsection*{Implementação (\texttt{quicksort.c})}

\begin{itemize}
  \item \textbf{Mediana-de-3} (\texttt{A[lo]}, \texttt{A[mid]},
    \texttt{A[hi]}) como pivô, evitando o pior caso $O(n^2)$ típico de
    entradas já ordenadas ou quase ordenadas.
  \item \textbf{Partição de Hoare}, com menos trocas em média que a
    partição de Lomuto.
  \item \textbf{Corte para inserção} em sub-arrays de até 16 elementos
    (\emph{insertion sort cutoff}): para $k$ pequeno, a constante menor
    do insertion sort compensa seu $O(k^2)$.
  \item \textbf{Eliminação de recursão de cauda}: recorre explicitamente
    apenas no lado \emph{menor} da partição, continuando em laço
    \texttt{while} no lado maior, limitando a pilha a $O(\log n)$ mesmo
    no pior caso.
\end{itemize}

\begin{lstlisting}[language=C]
static void quicksortRec(int *A, long int lo, long int hi) {
  while (hi - lo + 1 > INSERTION_CUTOFF) {
    long int mid = lo + (hi - lo) / 2;
    int pivot = medianOfThree(A, lo, mid, hi);
    long int p = hoarePartition(A, lo, hi, pivot);
    if (p - lo < hi - (p + 1)) {
      quicksortRec(A, lo, p);
      lo = p + 1;           /* continua no lado maior */
    } else {
      quicksortRec(A, p + 1, hi);
      hi = p;                /* continua no lado maior */
    }
  }
  insertionSort(A, lo, hi);
}
\end{lstlisting}

\begin{teorema}[Complexidade média do Quicksort]
O número esperado de comparações realizadas pelo Quicksort sobre um
array de $n$ elementos distintos, com pivôs tais que cada elemento tem
probabilidade $\Omega(1)$ de particionar razoavelmente bem, é
$\Theta(n \log n)$.
\end{teorema}
\begin{proof}[Esboço de demonstração (CLRS, cap.~7)]
Sejam $z_1 < z_2 < \dots < z_n$ os elementos em ordem. Para $i<j$, seja
$X_{ij}$ a variável indicadora de que $z_i$ e $z_j$ são comparados em
algum momento do algoritmo. Dois elementos só são comparados se um deles
for escolhido como pivô antes de qualquer elemento de valor
\emph{entre} eles; a probabilidade disso, num conjunto de $j-i+1$
elementos relevantes, é $\frac{2}{j-i+1}$. O número total esperado de
comparações é
\[
  E\!\left[\sum_{i=1}^{n-1}\sum_{j=i+1}^{n} X_{ij}\right]
  = \sum_{i=1}^{n-1}\sum_{j=i+1}^{n} \frac{2}{j-i+1} = O(n\log n),
\]
por linearidade da esperança e pela soma harmônica $\sum 1/k = O(\log n)$.
\end{proof}

\textbf{Pior caso:} $\Theta(n^2)$, quando a partição é maximamente
desbalanceada em toda chamada ($T(n) = T(n-1) + \Theta(n)$). A
mediana-de-3 torna esse caso raro na prática, mas não o elimina
adversarialmente. \textbf{Melhor caso:} $\Theta(n \log n)$, partições
balanceadas, $T(n) = 2T(n/2) + \Theta(n)$ (Teorema Mestre, caso~2).

\textbf{Complexidade espacial:} $O(\log n)$ esperado (pilha de
recursão, graças à eliminação de recursão de cauda); \textbf{sem}
alocação de nenhum array auxiliar de tamanho proporcional a $n$ ---
ordena \emph{in-place}.

\subsection{Radix Sort}

\textbf{Classe:} ordenação \textbf{não comparativa}, estável, por
distribuição (\emph{counting sort} repetido por dígito).

\textbf{Ideia central:} ordena-se repetidamente por "dígitos" do valor,
do menos significativo ao mais significativo (variante \textbf{LSD} ---
\emph{least-significant-digit-first}), usando em cada passada um
\emph{counting sort} estável.

\subsubsection*{Implementação (\texttt{radixsort.c})}

Base $256$ (dígitos de 8 bits) $\Rightarrow$ \textbf{4 passadas} para
cobrir os 32 bits de um \texttt{int}. Cada passada conta ocorrências de
cada um dos 256 valores do dígito corrente, converte em posições via
soma de prefixos, e distribui os elementos no array auxiliar; os papéis
de fonte/destino são trocados por ponteiro a cada passada.

\begin{lstlisting}[language=C]
static void countingPass(int *src, int *dst, long int n, int shift) {
  long int count[RADIX_BASE + 1];
  memset(count, 0, sizeof(count));
  for (long int i = 0; i < n; i++) {
    unsigned int digit = ((unsigned int)src[i] >> shift) & RADIX_MASK;
    count[digit + 1]++;
  }
  for (int d = 0; d < RADIX_BASE; d++) count[d + 1] += count[d];
  for (long int i = 0; i < n; i++) {
    unsigned int digit = ((unsigned int)src[i] >> shift) & RADIX_MASK;
    dst[count[digit]++] = src[i];
  }
}
\end{lstlisting}

\textbf{Complexidade temporal:} $O(d \cdot (n+b))$, com $d$ = número de
passadas, $b$ = base. Aqui $d=4$ e $b=256$ são \textbf{constantes
fixas} (decorrem só do tamanho fixo de 32 bits do \texttt{int}),
reduzindo a complexidade a $\Theta(n)$.

\begin{teorema}[Radix sort não viola $\Omega(n\log n)$]
O limite inferior $\Omega(n\log n)$ para ordenação vale apenas para
algoritmos que só podem comparar pares de chaves (prova pela árvore de
decisão: distinguir $n!$ permutações exige profundidade mínima
$\log_2(n!) = \Theta(n\log n)$). O radix sort nunca compara duas chaves
entre si --- usa o valor absoluto de cada dígito para indexar
diretamente posições no array, uma operação mais poderosa que
comparação binária, e por isso escapa legitimamente desse limite.
\end{teorema}

\textbf{Complexidade espacial: $O(n+b)$.} Diferente do quicksort, o
radix sort \textbf{não é in-place}: cada passada escreve num array
auxiliar de tamanho $n$. Em pico, mantém \textbf{dois arrays de tamanho
$n$} simultaneamente na memória --- o dobro do espaço de dados que o
quicksort usa. Essa característica se manifesta empiricamente na
Seção~\ref{sec:t2e9} como desvantagem prática para $n$ muito grande.

\subsection{Comparação teórica resumida}

\begin{table}[H]
\centering
\caption{Quicksort vs.\ Radix Sort --- comparação teórica.}
\begin{tabular}{l l l}
\toprule
\rowcolor{tablehead}
 & Quicksort (implementado) & Radix Sort (implementado)\\
\midrule
Classe & Comparação, \emph{in-place} & Distribuição, não-comparativo\\
Tempo médio & $\Theta(n\log n)$ & $\Theta(n)$ ($d,b$ fixos)\\
Pior caso & $\Theta(n^2)$ (raro, mitigado) & $\Theta(n)$ sempre\\
Espaço auxiliar & $O(\log n)$ (pilha) & $O(n+b)$ (array duplicado)\\
Estável & Não & Sim\\
\bottomrule
\end{tabular}
\end{table}

%======================================================================
\section{Metodologia e Implementação Computacional}
%======================================================================

\subsection{Ambiente de Execução}
\label{sec:ambiente}

\begin{table}[H]
\centering
\caption{Ambiente de hardware/software usado em todas as medições deste Relatório (computador pessoal do autor).}
{\small
\begin{tabular}{L{4.0cm} L{10.5cm}}
\toprule
\rowcolor{tablehead}
Item & Valor\\
\midrule
CPU & 13th Gen Intel(R) Core(TM) i7-1355U (1,70 GHz, 10 núcleos, 12 threads)\\
RAM total & 16.777.216 kB (16,0 GB)\\
RAM disponível (na coleta) & 16.454.352 kB\\
Memória virtual (commit) & 34.280.144 kB total, 28.839.404 kB em uso\\
Sistema operacional & Windows 11 Home Single Language, 25H2 (build 26200.9711, 64 bits)\\
Compilador & gcc (MinGW-W64 x86\_64-ucrt-posix-seh, build Brecht Sanders, r8) 13.2.0\\
Flags de compilação & \texttt{-O2}\\
\bottomrule
\end{tabular}
}
\end{table}

\subsection{Compilação}

\begin{lstlisting}
gcc -O2 rng.c utilarray.c gen_input.c -o gen_input.exe
gcc -O2 rng.c utilarray.c elapsedRunningTime.c quicksort.c radixsort.c run_sort.c -o run_sort.exe
gcc -O2 rng.c utilarray.c quicksort.c radixsort.c benchmark_clock.c -o benchmark_clock.exe
gcc -O2 rng.c secretkey_gen.c -o secretkey_gen.exe
\end{lstlisting}

\subsection{Medição do Tempo de Execução}
\label{sec:medicao}

O tempo é medido com \texttt{clock()} (tempo de CPU) e
\texttt{gettimeofday()} (relógio de parede), via
\texttt{elapsedRunningTime.c} fornecido pelo professor, em torno
\emph{apenas} da chamada ao algoritmo de ordenação (\texttt{start()} /
\texttt{stop()}), excluindo geração e leitura da entrada.

\textbf{Observação metodológica importante:} no Windows/MinGW usado,
\texttt{CLOCKS\_PER\_SEC = 1000}, ou seja, \texttt{clock()} tem
resolução de apenas $\sim$1\,ms --- grosseira demais para medir
corretamente execuções de microssegundos a poucos milissegundos (visível
nos dados brutos: \texttt{tempo\_clock\_s} aparece quantizado em
múltiplos de 0,001, enquanto \texttt{tempo\_wall\_s} tem resolução de
microssegundos). No Linux, \texttt{CLOCKS\_PER\_SEC} é tipicamente
1.000.000. Por isso, \textbf{o tempo de relógio de parede
(\texttt{gettimeofday()}) é usado como medida principal} nas análises
estatísticas deste Relatório (Seção~\ref{sec:resultados}); somente a
tabela da planilha oficial (Seção~\ref{sec:planilha-oficial}), que o
enunciado pede especificamente com \texttt{clock()}, usa esse relógio.

\subsection{Controle de Confusão Experimental}

Para cada tamanho $n$ e repetição, gera-se \textbf{um único array
aleatório} (\texttt{gen\_input}, semente $= n\cdot1000003 +
\mathit{repeticao}\cdot7919 + 12345$) e o \textbf{mesmo} array binário é
usado para testar os dois algoritmos, eliminando a variação de entrada
como fator de confusão e permitindo comparação pareada
(Seção~\ref{sec:estatistica}).

\subsection{Tamanhos e Repetições}

Tamanhos pequenos ($n\le10^4$) foram repetidos 7$\times$; médios
($10^5$--$10^6$) 3 a 5$\times$; grandes ($\ge5\times10^6$) 1 a
3$\times$, reduzindo repetições conforme o tempo por execução cresce.
Nenhuma execução atingiu o limite de 3600\,s estipulado pelo enunciado.
Os 14 tamanhos oficiais da planilha ($T(10^6)$ a $T(10^{10})$) foram
medidos à parte, numa única execução cada, por \texttt{benchmark\_clock.c}.

%======================================================================
\section{Resultados Experimentais e Discussão Científica}
\label{sec:resultados}
%======================================================================

\subsection{Tempo médio de execução}

\begin{table}[H]
\centering
\caption{Tempo médio de execução (relógio de parede), por tamanho de entrada. Fonte: \texttt{resultados\barra{}raw\_results.csv}, agregado por \texttt{resultados\barra{}graficos\_r\barra{}analise.R}.}
\label{tab:resultados}
\begin{tabular}{r r r r}
\toprule
\rowcolor{tablehead}
$n$ & Quicksort (s) & Radix Sort (s) & razão (qs/radix)\\
\midrule
1.000       & 0,0000377  & 0,0000057  & 6,6$\times$\\
10.000      & 0,000410   & 0,0000801  & 5,1$\times$\\
100.000     & 0,005271   & 0,000667   & 7,9$\times$\\
500.000     & 0,025271   & 0,003708   & 6,8$\times$\\
1.000.000   & 0,055610   & 0,008487   & 6,6$\times$\\
5.000.000   & 0,298064   & 0,037674   & 7,9$\times$\\
10.000.000  & 0,614553   & 0,078156   & 7,9$\times$\\
50.000.000  & 3,355480   & 0,377036   & 8,9$\times$\\
100.000.000 & 7,713330   & 0,951264   & 8,1$\times$\\
200.000.000 & 14,147100 (1 execução) & 1,563120 (1 execução) & 9,1$\times$\\
\bottomrule
\end{tabular}
\end{table}

A dispersão entre repetições (quanto cada tempo individual variou em
torno dessa média) não é mostrada célula a célula para manter a tabela
simples de ler; está disponível na Tabela~\ref{tab:desviopadrao}
(abaixo) para quem quiser o detalhe, e visualmente na
Figura~\ref{fig:boxplot}.

\begin{table}[H]
\centering
\caption{Desvio-padrão (dispersão entre repetições) correspondente à Tabela~\ref{tab:resultados}. Tamanhos com apenas 1 execução não têm desvio-padrão definido.}
\label{tab:desviopadrao}
\begin{tabular}{r r r}
\toprule
\rowcolor{tablehead}
$n$ & Quicksort (s) & Radix Sort (s)\\
\midrule
1.000       & 0,0000150 & 0,0000017\\
10.000      & 0,0000499 & 0,0000455\\
100.000     & 0,000451  & 0,000195\\
500.000     & 0,000871  & 0,000432\\
1.000.000   & 0,004236  & 0,001196\\
5.000.000   & 0,004433  & 0,001803\\
10.000.000  & 0,012138  & 0,002646\\
50.000.000  & 0,002232  & 0,019041\\
100.000.000 & 0,101421  & 0,078973\\
200.000.000 & --- (1 execução) & --- (1 execução)\\
\bottomrule
\end{tabular}
\end{table}

\figuraopcional{resultados/graficos_r/01_tempo_vs_n_loglog.png}{Tempo médio de execução vs.\ $n$, escala log-log, com barras de $\pm 1$ desvio-padrão.}{fig:tempon}{Fonte: \texttt{resultados\barra{}graficos\_r\barra{}analise.R}, dados em \texttt{resultados\barra{}raw\_results.csv}.}

\figuraopcional{resultados/graficos_r/02_razao_qs_radix.png}{Razão tempo(quicksort)/tempo(radix sort) por tamanho de entrada.}{fig:razao}{Fonte: \texttt{resultados\barra{}graficos\_r\barra{}analise.R}.}

\figuraopcional{resultados/graficos_r/03_boxplot_variabilidade.png}{Variabilidade entre repetições (tamanhos com $\ge$3 execuções).}{fig:boxplot}{Fonte: \texttt{resultados\barra{}graficos\_r\barra{}analise.R}.}

\subsection{Análise Estatística}
\label{sec:estatistica}

\textbf{Regressão log-log} ($\log(\mathit{tempo}) = \log(a) + k\log(n)$),
ajustada sobre as médias por tamanho, estima o expoente de crescimento
empírico:

\begin{table}[H]
\centering
\caption{Regressão log-log do tempo de execução (relógio de parede) em função de $n$.}
\begin{tabular}{l r r r}
\toprule
\rowcolor{tablehead}
Algoritmo & $k$ estimado & IC 95\% & $R^2$ ajustado\\
\midrule
Quicksort   & 1,0555 & [1,0459 ; 1,0651] & 0,9999\\
Radix Sort  & 1,0227 & [1,0018 ; 1,0436] & 0,9993\\
\bottomrule
\end{tabular}
\end{table}

Ambos os ajustes são excelentes ($R^2 > 0{,}999$). O expoente do
\textbf{quicksort (1,0555) é visivelmente maior} que o do \textbf{radix
sort (1,0227)}, refletindo o fator $\log n$ extra de
$\Theta(n\log n)$ presente no quicksort e ausente no radix sort. O
intervalo de confiança do radix sort \textbf{contém 1,0}, consistente
com $\Theta(n)$ teórico; o do quicksort \textbf{não contém 1,0},
evidência de crescimento genuinamente superlinear.

\textbf{Teste pareado} (Wilcoxon signed-rank, repetições com a mesma
entrada para os dois algoritmos como pares naturais, $n=68$ pares,
tamanhos de $10^3$ a $10^8$): $V=2346$, $p\approx7{,}8\times10^{-13}$. O
quicksort foi mais lento que o radix sort em \textbf{100\% dos 68 pares}
comparados --- diferença estatisticamente altamente significativa e
sistemática, não um artefato de média.

\subsection{Tamanhos extremos (planilha oficial, até 2 bilhões)}
\label{sec:planilha-oficial}
\label{sec:t2e9}

\begin{table}[H]
\centering
\caption{Tempos medidos com \texttt{clock()}, execução única, nos 14 tamanhos oficiais da planilha do professor. Fonte: \texttt{resultados\barra{}tempos\_planilha\_clock.csv}.}
\begin{tabular}{r r r}
\toprule
\rowcolor{tablehead}
Tamanho & Quicksort (s) & Radix Sort (s)\\
\midrule
$T(10^6)$  & 0,057   & 0,007\\
$T(2\times10^6)$  & 0,111   & 0,015\\
$T(5\times10^6)$  & 0,334   & 0,038\\
$T(10^7)$  & 0,618   & 0,079\\
$T(2\times10^7)$  & 1,286   & 0,156\\
$T(5\times10^7)$  & 3,415   & 0,414\\
$T(10^8)$  & 7,141   & 0,849\\
$T(2\times10^8)$  & 16,154  & 1,518\\
$T(5\times10^8)$  & 38,099  & 4,012\\
$T(10^9)$  & 81,306  & 12,322\\
$\mathbf{T(2\times10^9)}$  & \textbf{222,180} & \textbf{503,456}\\
$T(6\times10^9)$  & não testado & não testado\\
$T(8\times10^9)$  & não testado & não testado\\
$T(10^{10})$ & não testado & não testado\\
\bottomrule
\end{tabular}
\end{table}

\figuraopcional{resultados/graficos_r/04_planilha_oficial_loglog.png}{Tamanhos oficiais da planilha, até 2 bilhões de elementos.}{fig:planilha}{Fonte: \texttt{resultados\barra{}graficos\_r\barra{}analise.R}, dados em \texttt{resultados\barra{}tempos\_planilha\_clock.csv}.}

\textbf{O resultado em $T(2\times10^9)$ inverte a tendência}: pela
primeira vez o radix sort (503,5\,s) é \textbf{mais lento} que o
quicksort (222,2\,s). Não é erro de medição --- é consequência direta
da arquitetura de memória:

\begin{itemize}
  \item Em $n=2\times10^9$, um array de \texttt{int} (4 bytes) ocupa
    \textbf{8\,GB}.
  \item O \textbf{quicksort é in-place}: usa só esses 8\,GB.
  \item O \textbf{radix sort aloca um array auxiliar do mesmo tamanho}
    (Seção~2.2): em pico, precisa de $\sim$\textbf{16\,GB} ---
    praticamente toda a RAM física da máquina (16,0\,GB,
    Seção~\ref{sec:ambiente}).
  \item Isso força paginação em disco, ordens de grandeza mais lenta
    que acesso à RAM. O próprio ambiente registrado
    (Seção~\ref{sec:ambiente}) mostra \textbf{28,8\,GB de memória
    virtual em uso} num sistema com apenas 16\,GB de RAM física ---
    evidência direta de pressão de memória nesse momento.
  \item \textbf{Ressalva metodológica:} este ponto é uma única execução,
    sem repetição (diferente da Seção~\ref{sec:resultados}, que tem
    múltiplas repetições e desvio-padrão). A inversão é consistente com
    o raciocínio teórico acima e com a leitura de memória virtual
    registrada, mas não se pode descartar 100\% um efeito de ruído do
    sistema sem repetir a medição --- não feita para evitar o alto
    custo de tempo (cada execução de $T(2\times10^9)$ já leva vários
    minutos).
\end{itemize}

Este é um exemplo concreto de por que a \textbf{complexidade espacial
importa na prática, não só na teoria}: o radix sort tem melhor
complexidade de tempo assintótico ($\Theta(n)$ vs.\ $\Theta(n\log n)$),
mas sua maior exigência de memória ($O(n+b)$ com constante efetiva~2,
pelo array duplicado) o torna mais lento assim que o tamanho do
problema se aproxima do limite de RAM física disponível.

\subsection{Mesma planilha, medida com \texttt{gettimeofday()}}
\label{sec:gettimeofday}

A tabela da Seção~\ref{sec:planilha-oficial} usa \texttt{clock()},
que mede \textbf{tempo de CPU} (apenas os ciclos em que o processo
esteve de fato executando instruções). Os mesmos 14 tamanhos oficiais
foram remedidos com \texttt{gettimeofday()}, que mede \textbf{tempo de
relógio de parede} (\emph{wall-clock}): o intervalo real decorrido,
incluindo tempo em que o processo ficou bloqueado --- por exemplo,
aguardando o sistema operacional satisfazer uma falha de página durante
paginação em disco. Execução única por tamanho, como na medição
anterior.

\begin{table}[H]
\centering
\caption{Tempos medidos com \texttt{gettimeofday()} (tempo de relógio de
parede), execução única, nos mesmos 14 tamanhos oficiais da planilha do
professor. Fonte: \texttt{tempos\barra{}planilha\barra{}gettimeofday.csv}.}
\label{tab:gettimeofday}
\begin{tabular}{r r r}
\toprule
\rowcolor{tablehead}
Tamanho & Quicksort (s) & Radix Sort (s)\\
\midrule
$T(10^6)$  & 0,060103  & 0,009102\\
$T(2\times10^6)$  & 0,140042  & 0,016028\\
$T(5\times10^6)$  & 0,306008  & 0,045090\\
$T(10^7)$  & 0,672141  & 0,085645\\
$T(2\times10^7)$  & 1,330477  & 0,227832\\
$T(5\times10^7)$  & 4,320474  & 0,656880\\
$T(10^8)$  & 8,720561  & 1,087448\\
$T(2\times10^8)$  & 15,489256 & 1,646543\\
$T(5\times10^8)$  & 37,005234 & 9,020103\\
$T(10^9)$  & 110,744745 & 16,734537\\
$\mathbf{T(2\times10^9)}$  & \textbf{243,714830} & \textbf{632,344466}\\
$T(6\times10^9)$  & não testado & não testado\\
$T(8\times10^9)$  & não testado & não testado\\
$T(10^{10})$ & não testado & não testado\\
\bottomrule
\end{tabular}
\end{table}

\textbf{A inversão em $T(2\times10^9)$ se mantém --- e é ainda mais
acentuada}: sob relógio de parede, o radix sort (632,3\,s) é
\textbf{2,85$\times$ mais lento} que o quicksort (243,7\,s), contra
\textbf{2,27$\times$} na medição de tempo de CPU (503,5\,s vs.\
222,2\,s). Isso \textbf{reforça}, em vez de contradizer, a explicação
por pressão de memória da Seção anterior: o tempo gasto em troca de
páginas com o disco conta como tempo de relógio de parede decorrido,
mas conta pouco ou nada como tempo de CPU (o processo fica
\emph{bloqueado}, não executando instruções, enquanto o SO resolve a
falta de página) --- logo, é esperado que \texttt{gettimeofday()}
exponha esse custo de forma mais completa que \texttt{clock()}
justamente no ponto em que a paginação em disco é mais severa.

De forma mais geral, em todos os 11 tamanhos testados o tempo de
relógio de parede é igual ou maior que o tempo de CPU correspondente
(como deveria ser, já que nenhum processo pode gastar mais CPU do que
tempo decorrido), exceto por três pontos do quicksort
($T(5\times10^6)$, $T(2\times10^8)$ e $T(5\times10^8)$) em que o
\texttt{gettimeofday()} registrou um valor até 8\% \textbf{menor} que o
\texttt{clock()} --- fisicamente só possível como ruído de medição,
consistente com a mesma ressalva já feita na seção anterior: ambas as
medições são de execução única, sem repetição, e portanto sujeitas a
variação do sistema (outros processos, \emph{scheduling} do SO) que não
pode ser separada de um efeito real sem repetir a medição. O caso mais
chamativo é o radix sort em $T(5\times10^8)$: 4,012\,s por \texttt{clock()}
contra 9,020\,s por \texttt{gettimeofday()} (mais que o dobro) --- nesse
ponto o array auxiliar do radix sort já ocupa $\sim$4\,GB (2\,GB de
entrada + 2\,GB de cópia), o que começa a pressionar a memória bem
antes do colapso completo em $T(2\times10^9)$, mesmo sem ainda inverter
a ordem dos dois algoritmos.

\subsection{Tamanhos não testados ($T(6\times10^9)$, $T(8\times10^9)$, $T(10^{10})$)}

Não testados por \textbf{dois limites independentes e cumulativos}:

\begin{enumerate}
  \item \textbf{RAM física}: 6, 8 e 10 bilhões de \texttt{int}s exigiriam
    $\sim$24, 32 e 40\,GB, respectivamente --- mais que os 16\,GB
    disponíveis, mesmo antes do array auxiliar do radix sort.
  \item \textbf{Modelo de dados do Windows (LLP64)}: o campo
    \texttt{long int length} do \texttt{struct Array} (fornecido pelo
    professor) tem \textbf{64 bits no Linux} (modelo LP64, usado nos
    exemplos originais do professor), mas apenas \textbf{32 bits no
    Windows} (modelo LLP64, adotado pelo MinGW-w64 por seguir a ABI do
    Windows) --- mesmo em compilação de 64 bits. Isso limita o tamanho
    máximo representável a $2^{31}-1 \approx 2{,}147$ bilhões,
    independentemente da RAM disponível. É uma diferença de
    portabilidade real entre Linux e Windows, documentada aqui, não uma
    limitação do algoritmo.
\end{enumerate}

Essas células foram marcadas "Sem memória" por uma verificação prévia
no código (\texttt{benchmark\_clock.c}), que evita a tentativa de
alocação, já que seria inútil (RAM) e, para $n > 2^{31}-1$, produziria
comportamento indefinido pelo \emph{overflow} do campo \texttt{length}.

\subsection{Conclusão}

Os dados confirmam a teoria (Seção~2) em praticamente toda a faixa
testada: o radix sort é consistentemente de 5 a 9$\times$ mais rápido
que o quicksort, com vantagem relativa que \textbf{cresce} com $n$ (de
$\sim$6,6$\times$ em $n=10^3$ a $\sim$9,1$\times$ em $n=2\times10^8$)
--- compatível com o fator $\log n$ extra do quicksort, que pesa cada
vez mais conforme $n$ cresce. A regressão log-log confirma
estatisticamente a diferença de expoentes e o teste pareado confirma
que a vantagem do radix sort é sistemática ($p<10^{-12}$). Porém, em
$n=2\times10^9$, essa vantagem se inverte na prática por causa do maior
custo de memória do radix sort colidir com o limite físico de RAM da
máquina --- um lembrete de que complexidade assintótica de tempo e
desempenho real divergem quando a complexidade \textbf{espacial} de um
algoritmo o aproxima dos limites físicos do hardware disponível.

%======================================================================
\section{Questões do Questionário do EP2}
%======================================================================

\subsection{Questão 1 --- Chave secreta}

Método idêntico ao de \texttt{calcSecret.c}: inicializa o gerador
congruente linear (constantes de \texttt{java.util.Random}:
multiplicador \texttt{0x5DEECE66D}, incremento \texttt{0xB}, módulo
$2^{48}$) com \texttt{setSeed(semente)} e aplica \texttt{next(31)}
repetidamente; o último valor retornado é a chave.

\begin{itemize}
  \item Semente: RA do autor = \texttt{18245598}
  \item Iterações: \texttt{31}
  \item \textbf{Chave secreta: \texttt{1906028835}}
\end{itemize}

Valor puramente aritmético (independe de hardware/SO/compilador),
reconfirmado executando \texttt{secretkey\_gen.exe 18245598 31} no
mesmo ambiente Windows usado para os benchmarks.

\subsection{Questão 2 --- Crítica ao método de geração da chave}

O gerador (LCG --- \emph{Linear Congruential Generator}) é o mesmo de
\texttt{java.util.Random}: $\mathit{seed} = (\mathit{seed}\cdot
\texttt{0x5DEECE66D} + \texttt{0xB}) \bmod 2^{48}$,
$\mathit{next}(\mathit{bits}) = \mathit{seed} \gg (48-\mathit{bits})$.

\textbf{Pontos fortes para este uso:} determinístico e reprodutível
(mesma semente $\Rightarrow$ mesma chave, sem precisar armazenar
tabela); rápido; para fins de \emph{identificação} (não de segurança), a
qualidade estatística do LCG é irrelevante.

\textbf{Fraquezas (crítica técnica, mesmo não comprometendo este uso):}

\begin{itemize}
  \item LCGs têm \textbf{estrutura em reticulado} (Marsaglia, 1968,
    "Random numbers fall mainly in the planes"): $k$-uplas consecutivas
    caem num número reduzido de hiperplanos paralelos, nunca
    preenchendo o espaço uniformemente.
  \item Os \textbf{bits baixos de um LCG têm período curto} --- por isso
    \texttt{next(bits)} usa só os bits altos do estado. O uso de
    \texttt{next(31)} (31 bits mais significativos) é a escolha correta
    sob esse aspecto.
  \item Dado o módulo $2^{48}$ e 31 iterações fixas, há mais valores de
    saída possíveis (até $\sim2\times10^9$) do que RAs plausíveis,
    tornando colisões entre dois RAs diferentes improváveis mas não
    impossíveis (paradoxo do aniversário) em turmas de dezenas a poucas
    centenas de alunos.
  \item Um LCG \textbf{não é adequado para chaves criptográficas}
    (previsível a partir de poucas saídas consecutivas) --- irrelevante
    aqui, que é uso de identificação pública, não de sigilo.
\end{itemize}

\subsection{Questão 3 --- Código-fonte}

Entregue em \texttt{src/}: \texttt{rng.c\barra{}h},
\texttt{utilarray.c\barra{}h}, \texttt{elapsedRunningTime.c\barra{}h}
(fornecidos pelo professor, sem alterações); \texttt{quicksort.c\barra{}h},
\texttt{radixsort.c\barra{}h} (implementação própria, Seção~2);
\texttt{gen\_input.c}, \texttt{run\_sort.c}, \texttt{benchmark\_clock.c},
\texttt{secretkey\_gen.c} (programas de suporte, próprios). Compilação
na Seção~3.2.

\subsection{Questão 4 --- Preenchimento da planilha}

Os dados necessários para o preenchimento --- identificador
(\texttt{1906028835}), informações de ambiente (Seção~\ref{sec:ambiente})
e os 14 tempos oficiais por algoritmo (Tabela em
Seção~\ref{sec:planilha-oficial}, disponíveis também em
\texttt{resultados\barra{}tempos\_planilha\_clock.csv}) --- estão prontos
e documentados neste Relatório.

\subsection{Questão 5 --- Discussão científica}

Ver Seção~\ref{sec:resultados} (Resultados Experimentais e Discussão
Científica), que cobre metodologia, resultados tabulados, análise
estatística (regressão log-log e teste pareado de Wilcoxon), o achado do
ponto de inversão em $T(2\times10^9)$ e sua explicação em termos de
complexidade espacial, e os limites de representação em
\texttt{long int} no Windows.

%======================================================================
\section{Fontes e Referências}
%======================================================================

\begin{itemize}
  \item F.~Nakano. \emph{Enunciado e código-base do EP2, IAA}.
    \urlRepoEP.
  \item F.~Nakano. \emph{getPlatformInfo}.
    \urlPlatformInfo.
  \item T.~H.~Cormen, C.~E.~Leiserson, R.~L.~Rivest, C.~Stein.
    \emph{Introduction to Algorithms}, 3ª ed.\ --- análise do quicksort
    por variáveis indicadoras (cap.~7) e do radix sort (cap.~8).
  \item G.~Marsaglia. ``Random numbers fall mainly in the planes''.
    \emph{Proceedings of the National Academy of Sciences}, 61(1), 1968
    --- estrutura em reticulado de LCGs.
  \item Oracle. Documentação de \texttt{java.util.Random} --- constantes
    do gerador congruente linear usado em \texttt{rng.c}.
  \item R Core Team. \emph{R: A Language and Environment for Statistical
    Computing}, versão 4.3.3 --- usado em
    \texttt{resultados\barra{}graficos\_r\barra{}analise.R} para a
    regressão log-log, teste de Wilcoxon e os gráficos deste Relatório.
\end{itemize}

\end{document}
